% Bài báo #1 - bản thảo tiếng Việt (cùng nội dung, số liệu và trích dẫn với main_en.tex).
% Biên dịch: tectonic main_vi.tex   (XeTeX + BibTeX, xem README.md)
\documentclass[a4paper,11pt]{article}

\usepackage{fontspec}
\setmainfont{texgyretermes}[Extension=.otf,UprightFont=*-regular,BoldFont=*-bold,ItalicFont=*-italic,BoldItalicFont=*-bolditalic]
\usepackage[english,vietnamese]{babel}
\setlength{\emergencystretch}{2em}
\usepackage{amsmath,amssymb}
\usepackage{icomma} % dấu phẩy thập phân trong công thức: 0,874
\usepackage{graphicx}
\usepackage{booktabs}
\usepackage{multirow}
\usepackage[margin=2.4cm]{geometry}
\usepackage[numbers,sort&compress]{natbib}
\usepackage{caption}
\captionsetup{font=small,labelfont=bf}
\usepackage{tikz}
\usetikzlibrary{arrows.meta,positioning,fit}
\usepackage[colorlinks=true,linkcolor=blue!50!black,citecolor=blue!50!black,urlcolor=blue!50!black]{hyperref}

\newcommand{\nuot}{\nu_{12}}
\newcommand{\nuto}{\nu_{21}}
\newcommand{\RR}{R^2}
% babel đặt lại caption khi vào document -> ghi đè qua \captionsvietnamese.
\addto\captionsvietnamese{%
  \renewcommand{\refname}{Tài liệu tham khảo}%
  \renewcommand{\figurename}{Hình}%
  \renewcommand{\tablename}{Bảng}%
  \renewcommand{\abstractname}{Tóm tắt}}

\title{Thiết kế ngược vật liệu meta auxetic bằng mô hình sinh sâu: cân bằng
yêu cầu đa mục tiêu và khả năng tổng quát hóa ngoài phân phối}
\author{Trịnh Bình Minh\thanks{[Đơn vị công tác và thông tin tác giả liên hệ - sẽ bổ sung.]}}
\date{Bản thảo, 30/09/2026}

\begin{document}
\maketitle

\begin{abstract}
Các mô hình sinh sâu có thể đề xuất ô cơ sở vật liệu meta cho một tính chất
hiệu dụng mục tiêu chỉ trong vài mili giây. Tuy nhiên, một thiết kế dùng được
phải đồng thời đạt ba yêu cầu. Nó phải khớp mục tiêu sau khi nhị phân hóa và
kiểm chứng bằng phần tử hữu hạn (FE). Nó phải chế tạo được. Và lý tưởng là
mô hình phải tạo được thiết kế cho cả những mục tiêu nằm ngoài dữ liệu huấn
luyện. Chúng tôi nghiên cứu ba yêu cầu này cho ô cơ sở auxetic hai chiều
được sinh bởi một bộ mã hóa tự động biến phân có điều kiện (cVAE), tinh
chỉnh bằng đồng nhất hóa khả vi. Về độ chính xác, chúng tôi chỉ ra rằng tinh
chỉnh không gian ẩn lúc suy luận phải tối ưu chính thiết kế được kiểm chứng.
Tinh chỉnh trên đầu ra liên tục của bộ giải mã đưa hàm mục tiêu về
$\sim\!10^{-10}$, trong khi sai số kiểm chứng vẫn lớn hơn sáu bậc. Ngược lại,
khi đưa các phép biến đổi của khâu kiểm chứng vào hàm mục tiêu, $\RR$ kiểm
chứng FE của $\nuot$ trên 100 mục tiêu kiểm tra tăng từ 0,874 lên 0,998 (một
mẫu) và từ 0,986 lên 0,997 (tốt nhất trong 30 mẫu). Về chọn lọc đa mục
tiêu, điểm tổng hợp có trọng số giữa độ chính xác, khả năng chế tạo và thẩm
mỹ tương quan ở mức vừa phải với một phép xếp hạng Pareto độc lập (trung vị
Spearman 0,71). Với mục tiêu ngoài phân phối, cVAE sinh được thiết kế có
dấu hệ số Poisson đúng nhưng chưa từng xuất hiện trong dữ liệu, điều mà truy
xuất láng giềng gần nhất không làm được. Tuy vậy, mô hình không ngoại suy
được biên độ ra ngoài dải huấn luyện.
\end{abstract}

\noindent\textbf{Từ khóa:} Vật liệu meta auxetic; Mô hình sinh sâu; Vật lý
khả vi; Thiết kế ngược; Tối ưu hóa tô-pô; Đồng nhất hóa

% ---------------------------------------------------------------------------
\section{Mở đầu}

Vật liệu meta auxetic co lại theo phương ngang khi bị nén một trục. Hệ số
Poisson âm của chúng bắt nguồn từ vi kiến trúc chứ không từ thành phần hóa
học của vật liệu cấu thành
\citep{lakes1987foam,evans2000auxetic,bertoldi2017flexible}. Đồng nhất hóa
ngược bằng tối ưu hóa tô-pô là con đường kinh điển để đi từ một tensor hiệu
dụng cho trước đến ô cơ sở
\citep{sigmund1994materials,larsen1997design,bendsoe2003topology}, nhưng mỗi
mục tiêu mới đòi hỏi một bài toán tối ưu lặp với hàng trăm lần giải FE.
Thiết kế ngược dựa trên dữ liệu phân bổ chi phí này ra trước: bộ mã hóa tự
động biến phân, mạng đối sinh và mô hình khuếch tán, khi đã huấn luyện trên
thư viện mô phỏng, trả về kiến trúc ứng viên gần như tức thì
\citep{wang2020deep,mao2020gan,kumar2020spinodoid,bastek2022inverting,kollmann2020deep,zheng2021controllable,pahlavani2024deep,bastek2023diffusion}.

Khi dùng thực tế, một mô hình sinh như vậy phải đáp ứng ba yêu cầu, nhưng
chúng hiếm khi được xem xét cùng nhau.
\emph{(i) Độ chính xác đã kiểm chứng.} Thứ được chế tạo là phiên bản đã nhị
phân hóa và lấy mẫu lại của trường mật độ sinh ra, nên độ chính xác phải
được đánh giá bằng phân tích FE trên chính thiết kế đó, không phải bằng mô
hình thay thế (surrogate) học được. Surrogate của ánh xạ thuận đã được biết
là có thể bị bộ tối ưu và mô hình sinh khai thác \citep{trabucco2021coms}.
Một mẫu đơn lẻ cũng hiếm khi đạt dung sai chặt, nên các quy trình phải sinh
nhiều ứng viên rồi chọn bằng FE \citep{pahlavani2024deep}.
\emph{(ii) Đa mục tiêu.} Một thiết kế khớp mục tiêu nhưng có đảo vật liệu
rời rạc, thanh mảnh dưới độ phân giải chế tạo, hoặc biên không ghép lát được
thì không dùng được. Do đó độ chính xác phải được cân bằng với khả năng chế
tạo, và ở mức ưu tiên thấp hơn, với tính đều đặn hình học.
\emph{(iii) Mục tiêu ngoài phân phối (OOD).} Lý do chính để dùng mô hình sinh
thay vì tra cứu trong thư viện huấn luyện là nó có thể tạo thiết kế cho
những mục tiêu mà thư viện không chứa. Việc nó có làm được hay không, và
theo hướng nào trong không gian tính chất, là câu hỏi thực nghiệm.

Bài báo này giải quyết ba yêu cầu trên bằng một quy trình và một giao thức
đánh giá duy nhất. Trong đó mọi con số đều được kiểm chứng FE trên thiết kế
đã nhị phân hóa, và mọi phép so sánh đều kèm khoảng tin cậy ghép cặp. Các
đóng góp của chúng tôi gồm:
\begin{enumerate}
\item \textbf{Tinh chỉnh không gian ẩn nhận biết nhị phân hóa.} Chúng tôi
tinh chỉnh vector ẩn của một cVAE cố định bằng gradient giải tích chính xác
của bộ giải đồng nhất hóa tuần hoàn. Chúng tôi chỉ ra rằng tinh chỉnh chỉ
đáng tin khi hàm mục tiêu được tính trên thiết kế được kiểm chứng, tức sau
khi ép biên tuần hoàn, chiếu Heaviside với tham số tăng dần
\citep{guest2004achieving,wang2011projection} và lấy mẫu lại láng giềng gần
nhất. Chúng tôi cũng đưa ra bằng chứng trực tiếp về hiện tượng khai thác
vật liệu xám xảy ra nếu không làm vậy.
\item \textbf{Chọn lọc đa mục tiêu có kiểm tra độc lập.} Chúng tôi xếp hạng
ứng viên bằng điểm tổng hợp có trọng số giữa độ chính xác, khả năng chế tạo
và thẩm mỹ, rồi kiểm tra điểm này với một phép xếp hạng Pareto độc lập thay
vì mặc định nó hợp lý.
\item \textbf{Đặc trưng hóa OOD có cơ sở vật lý.} Chúng tôi chỉ ra rằng các
mục tiêu đối xứng $\nuot=\nuto\le-1$ không chấp nhận được với ô cơ sở trực
hướng, khiến các bộ đánh giá ``cực trị'' ngây thơ trở nên vô hiệu. Sau đó
chúng tôi tách khả năng tổng quát hóa theo \emph{dấu} của hệ số Poisson, nơi
cVAE thành công còn truy xuất thất bại, khỏi khả năng tổng quát hóa theo
\emph{biên độ}, nơi mọi phương pháp đều bão hòa ở biên dữ liệu huấn luyện.
\item \textbf{Kết quả âm tính có kiểm soát.} Cả đầu hồi quy
Kolmogorov--Arnold (KAN) lẫn bộ giải mã biểu diễn ẩn wavelet (WIRE) đều
không cải thiện so với baseline tích chập khi đo bằng chỉ số kiểm chứng FE.
\end{enumerate}

% ---------------------------------------------------------------------------
\section{Kết quả}

\subsection{Khung phương pháp và giao thức đánh giá}

Hình~\ref{fig:framework} trình bày quy trình. cVAE ánh xạ cặp mục tiêu
$(\nuot^{*},\nuto^{*})$ cùng vector ẩn $z\sim\mathcal N(0,I)$ thành ảnh mật
độ $64\times64$. cVAE được huấn luyện hai giai đoạn; giai đoạn hai dùng loss
tính bằng đồng nhất hóa FE trên các mẫu tiên nghiệm (mục Phương pháp). Với
mỗi mục tiêu, $N$ vector ẩn được giải mã và hậu xử lý. Hậu xử lý gồm ép biên
tuần hoàn, lấy ngưỡng 0,5 và lấy mẫu lại láng giềng gần nhất về lưới FE
$50\times50$. Mỗi ứng viên sau đó được đồng nhất hóa FE và chấm điểm khả
năng chế tạo cùng thẩm mỹ. Ứng viên tốt nhất theo điểm tổng hợp có thể được
tinh chỉnh tiếp trong không gian ẩn. Thiết kế đã tinh chỉnh chỉ được giữ nếu
sai số kiểm chứng của nó thấp hơn (tinh chỉnh ``có kiểm soát'').

Mọi số liệu độ chính xác đều được kiểm chứng FE trên thiết kế đã nhị phân
hóa. Bộ đánh giá trong phân phối (IN100) gồm 100 điều kiện của tập kiểm tra,
rút với seed cố định. Các phép so sánh kèm khoảng tin cậy (KTC) 95\%
bootstrap ghép cặp theo điều kiện \citep{efron1993bootstrap}. Tiêu chí thành
công được ghi lại trước mỗi thí nghiệm (mục Phương pháp).

\begin{figure}[t]
\centering
\begin{tikzpicture}[
  node distance=3mm and 4mm,
  box/.style={draw, rounded corners=2pt, align=center, font=\small,
              minimum height=10mm, inner sep=3pt, fill=#1},
  box/.default=white,
  arr/.style={-{Stealth[length=2mm]}, thick}]
\node[box=gray!10] (t) {mục tiêu\\$(\nu_{12}^{*},\nu_{21}^{*})$\\$z_{1..N}\!\sim\!\mathcal N(0,I)$};
\node[box=orange!15, right=of t] (g) {bộ giải mã\\cVAE (tinh chỉnh\\vật lý)};
\node[box=blue!8, right=of g] (v) {ép biên tuần hoàn\\ngưỡng 0,5\\lấy mẫu $50^2$};
\node[box=green!12, right=of v] (fe) {đồng nhất\\hóa FE\\$\nu_{12},\nu_{21}$};
\node[box=yellow!15, right=of fe] (s) {điểm tổng hợp\\0,6 chính xác\\0,3 chế tạo\\0,1 thẩm mỹ};
\node[box=blue!15, right=of s] (r) {tinh chỉnh\\ẩn có\\kiểm soát};
\draw[arr] (t) -- (g);
\draw[arr] (g) -- (v);
\draw[arr] (v) -- (fe);
\draw[arr] (fe) -- (s);
\draw[arr] (s) -- (r);
\draw[arr, dashed] (r.north) -- ++(0,5mm) -| node[pos=0.25, above, font=\scriptsize]{kiểm chứng lại thiết kế đã tinh chỉnh} (fe.north);
\end{tikzpicture}
\caption{Quy trình thiết kế ngược. Mọi ứng viên đều được kiểm chứng bằng
đồng nhất hóa FE sau đúng các bước hậu xử lý mà một thiết kế đem chế tạo sẽ
trải qua. Bước chọn kết hợp độ chính xác đã kiểm chứng với khả năng chế tạo
(liên thông, kích thước chi tiết tối thiểu, tuần hoàn) và điểm thẩm mỹ trọng
số thấp. Chi tiết bước tinh chỉnh ở Hình~\ref{fig:pipeline}.}
\label{fig:framework}
\end{figure}

\subsection{Độ chính xác đã kiểm chứng đòi hỏi tối ưu chính thiết kế được kiểm chứng}

\paragraph{Tinh chỉnh liên tục khai thác vật liệu xám.}
Tinh chỉnh đầu ra liên tục của bộ giải mã (30 bước L-BFGS
\citep{liu1989lbfgs}, gradient qua bộ giải mã và qua độ nhạy chính xác của
đồng nhất hóa) có ích cho mẫu đơn lẻ của mô hình sản xuất. Trên IN100,
$\RR(\nuot)$ tăng từ 0,874 lên 0,976 và sai số tuyệt đối trung bình (MAE)
của $\nuot$ giảm 57,7\% (KTC 95\%: 49,4--64,6\%; Bảng~\ref{tab:main},
Hình~\ref{fig:parity}b). Tuy nhiên, cùng phép tinh chỉnh đó lại \emph{làm
xấu đi} ứng viên tốt nhất trong 30 mẫu. MAE($\nuot$) tăng từ 0,0152 lên
0,0246, tức thay đổi $-62{,}5\%$ (KTC $-98{,}3$ đến $-33{,}9\%$), và $\RR$
giảm từ 0,986 xuống 0,968. Diễn biến hàm mục tiêu giải thích nguyên nhân
(Hình~\ref{fig:gap}). Khi kết thúc tinh chỉnh liên tục, trung vị hàm mục
tiêu là $2{,}5\times10^{-10}$, nhưng trung vị bình phương sai số kiểm chứng
của thiết kế đã nhị phân hóa là $4{,}4\times10^{-4}$. Bộ tối ưu khớp mục
tiêu trên một trường có mật độ trung gian không còn tồn tại sau khi lấy
ngưỡng. Đây chính là dạng thất bại của việc khai thác surrogate, nhưng xuất
hiện ngay bên trong một mô hình vật lý chính xác.

\paragraph{Tinh chỉnh nhận biết nhị phân hóa.}
Khi đặt các phiên bản khả vi của phép biến đổi kiểm chứng vào trong hàm mục
tiêu (Hình~\ref{fig:pipeline}; 32 bước trên $\beta\in\{1;4;16;64\}$), khoảng
cách giữa hàm mục tiêu và kết quả kiểm chứng giảm ba bậc
(Hình~\ref{fig:gap}). Mọi tiêu chí trong phân phối đặt trước đều đạt
(Bảng~\ref{tab:main}).
\begin{itemize}
\item \textbf{Mẫu đơn lẻ.} $\RR(\nuot)$ tăng từ 0,874 lên 0,998 và
MAE($\nuot$) giảm 91,8\% (KTC 89,0--93,9\%). Sai số tương đối trung bình của
$\nuot$ giảm từ 14,4\% xuống 2,6\%, và $\RR(\nuto)$ tăng từ 0,477 lên 0,883.
\item \textbf{Sau khi chọn tốt nhất trong 30 mẫu.} MAE($\nuot$) giảm 62,6\%
(KTC 50,7--72,1\%) và $\RR(\nuot)$ tăng từ 0,986 lên 0,997. Tinh chỉnh có
kiểm soát chấp nhận 92\% số lần tinh chỉnh và giảm MAE($\nuot$) 69,1\% (KTC
61,6--75,4\%), không tốn thêm lần giải FE nào.
\end{itemize}
Sau bước chọn tốt nhất trong 30 mẫu, MAE($\nuto$) giảm từ 0,0327 xuống
0,0128, trong khi $\RR(\nuto)$ giảm từ 0,802 xuống 0,712. Như vậy một số điều
kiện có sai lệch $\nuto$ lớn hơn, và quy tắc kiểm soát dựa trên tổng sai số
không phải lúc nào cũng loại được chúng. Tinh chỉnh tốn trung bình 77 lần
giải FE mỗi điều kiện. Với 50--100\,ms mỗi lần giải kèm độ nhạy trên lưới
$50\times50$, chi phí này tương đương vài giây cho mỗi thiết kế.

\begin{figure}[t]
\centering
\begin{tikzpicture}[
  node distance=4mm and 5mm,
  box/.style={draw, rounded corners=2pt, align=center, font=\small,
              minimum height=9mm, inner sep=3pt, fill=#1},
  box/.default=white,
  arr/.style={-{Stealth[length=2mm]}, thick},
  grad/.style={-{Stealth[length=2mm]}, thick, dashed, blue!60!black}]
\node[box=gray!10] (z) {vector ẩn $z$\\(khởi tạo)};
\node[box=gray!10, below=of z] (c) {mục tiêu\\$(\nu_{12}^{*},\nu_{21}^{*})$};
\node[box=orange!15, right=8mm of z, yshift=-6.5mm] (dec) {bộ giải mã\\cVAE (cố định)};
\node[box=blue!8, right=of dec] (per) {ép biên\\tuần hoàn};
\node[box=blue!8, right=of per] (hv) {Heaviside\\$H_\beta(\rho)$};
\node[box=blue!8, right=of hv] (rs) {gần nhất\\$64^2\!\to\!50^2$};
\node[box=green!12, right=of rs] (fe) {đồng nhất hóa FE\\$\mathbf{Q}\to\nu_{12},\nu_{21}$};
\node[box=gray!10, below=9mm of hv] (loss) {$\mathcal{L}=\tfrac12\lVert\boldsymbol{\nu}-\boldsymbol{\nu}^{*}\rVert^2$};
\draw[arr] (z) -| ([xshift=-3mm]dec.west) -- (dec.west);
\draw[arr] (c) -| ([xshift=-3mm]dec.west);
\draw[arr] (dec) -- (per);
\draw[arr] (per) -- (hv);
\draw[arr] (hv) -- (rs);
\draw[arr] (rs) -- (fe);
\draw[arr] (fe.south) |- (loss.east);
\draw[grad] (loss.west) -- ([xshift=-6mm]loss.west -| z.west)
  node[pos=0.45, below, font=\scriptsize]{$\partial\mathcal{L}/\partial z$ chính xác, L-BFGS; $\beta\in\{1;4;16;64\}$}
  |- (z.west);
\node[draw=blue!50, dashed, rounded corners, fit=(per)(hv)(rs), inner sep=2pt,
      label={[font=\scriptsize]above:cùng các phép biến đổi như khâu kiểm chứng}] {};
\end{tikzpicture}
\caption{Tinh chỉnh không gian ẩn nhận biết nhị phân hóa. Đầu ra của bộ giải
mã cố định đi qua các phiên bản khả vi của những phép biến đổi dùng khi kiểm
chứng rồi mới vào đồng nhất hóa FE, nên đại lượng được tối ưu chính là đại
lượng được kiểm chứng. Biến thể liên tục bỏ khối nét đứt và dùng nội suy
song tuyến tính. Khâu kiểm chứng thay $H_\beta$ bằng ngưỡng cứng 0,5.}
\label{fig:pipeline}
\end{figure}

\begin{table}[t]
\centering
\caption{Độ chính xác kiểm chứng FE trên IN100 (100 mục tiêu kiểm tra, cVAE
sản xuất). $\Delta$MAE là mức giảm tương đối MAE($\nuot$) so với thiết kế
chưa tinh chỉnh cùng chế độ, kèm KTC 95\% bootstrap ghép cặp (10\,000 lần lấy
mẫu lại). Giá trị âm nghĩa là tinh chỉnh làm kết quả xấu đi. Số lần gọi FE
tính trên mỗi điều kiện, gồm cả bước chọn.}
\label{tab:main}
\small
\setlength{\tabcolsep}{4pt}
\begin{tabular}{llccccr}
\toprule
Chế độ & Tinh chỉnh & $\RR(\nuot)$ & MAE($\nuot$) & $\RR(\nuto)$ & $\Delta$MAE [KTC 95\%] & Lần gọi FE\\
\midrule
\multirow{3}{*}{Mẫu đơn lẻ}
 & không                  & 0,874 & 0,0517 & 0,477 & -- & 1\\
 & liên tục               & 0,976 & 0,0219 & 0,643 & $+57{,}7\%$ [49,4; 64,6] & 51\\
 & nhận biết nhị phân hóa & \textbf{0,998} & \textbf{0,0043} & \textbf{0,883} & $\mathbf{+91{,}8\%}$ [89,0; 93,9] & 78\\
\midrule
\multirow{4}{*}{Tốt nhất / 30}
 & không                  & 0,986 & 0,0152 & 0,802 & -- & 30\\
 & liên tục               & 0,968 & 0,0246 & 0,637 & $-62{,}5\%$ [$-98{,}3$; $-33{,}9$] & 79\\
 & nhận biết nhị phân hóa & 0,997 & 0,0057 & 0,712 & $+62{,}6\%$ [50,7; 72,1] & 107\\
 & \quad có kiểm soát     & \textbf{0,998} & \textbf{0,0047} & -- & $\mathbf{+69{,}1\%}$ [61,6; 75,4] & 107\\
\bottomrule
\end{tabular}
\end{table}

\begin{figure}[t]
\centering
\includegraphics[width=\linewidth]{figures/fig_parity_vi.pdf}
\caption{$\nuot$ kiểm chứng FE so với mục tiêu trên IN100 ($n=100$, mẫu đơn
lẻ của cVAE sản xuất): (a) chưa tinh chỉnh; (b) sau tinh chỉnh liên tục;
(c) sau tinh chỉnh nhận biết nhị phân hóa. Ở cả ba ô, mỗi điều kiện xuất
phát từ cùng một vector ẩn. Nét đứt: đường đồng nhất.}
\label{fig:parity}
\end{figure}

\begin{figure}[t]
\centering
\includegraphics[width=0.55\linewidth]{figures/fig_gap_vi.pdf}
\caption{Khoảng cách giữa hàm mục tiêu và kết quả kiểm chứng. Mỗi điểm là
một điều kiện của IN100. Trục ngang là giá trị cuối của hàm mục tiêu tinh
chỉnh; trục dọc là bình phương sai số hệ số Poisson của thiết kế đã nhị phân
hóa, đo bằng một lần giải FE độc lập. Tinh chỉnh liên tục (hình tròn) đạt
hàm mục tiêu gần $10^{-10}$ trong khi sai số kiểm chứng vẫn ở mức
$10^{-4}$--$10^{-2}$. Tinh chỉnh nhận biết nhị phân hóa (hình vuông) kéo các
điểm về gần đường đồng nhất. Giá trị dưới $10^{-14}$ được cắt để hiển thị.}
\label{fig:gap}
\end{figure}

\subsection{Yêu cầu đa mục tiêu: khả năng chế tạo và xếp hạng tổng hợp}

Thiết kế chính xác chưa chắc đã chế tạo được. Mỗi ứng viên được kiểm tra
(a) tính liên thông, không có thanh mảnh hơn hai pixel, và (b) tính tuần
hoàn, tức các cạnh đối diện phải khớp để ô cơ sở ghép lát được bằng phép
tịnh tiến (mục Phương pháp). Với một cVAE trước đây chưa ép biên tuần hoàn,
chỉ 0--3,5\% mẫu sinh ra qua được cả hai kiểm tra. Tỉ lệ này đồng đều trên
toàn không gian tính chất, nên vấn đề nằm ở hành vi của bộ giải mã chứ không
phải ở độ phủ dữ liệu. Ép tuần hoàn bằng cách lấy trung bình các cạnh đối
diện trước khi nhị phân hóa là một phép toán không cần huấn luyện. Nó đã
nâng tỉ lệ qua kiểm tra lên hơn một bậc, và được áp dụng trong mọi lần kiểm
chứng ở bài báo này. Với phép ép này, 24,7\% mẫu của mô hình sản xuất qua
được mọi kiểm tra (300 mục tiêu kiểm tra, $N=30$; Bảng~\ref{tab:multi}).
Trong phân phối, phép ép tuần hoàn vì vậy không buộc phải đánh đổi với độ
chính xác.

Trong số các ứng viên của một mục tiêu, bước chọn cực đại hóa điểm tổng hợp
$0{,}6\,s_{\mathrm{acc}}+0{,}3\,s_{\mathrm{man}}+0{,}1\,s_{\mathrm{aes}}$.
Các trọng số mã hóa thứ tự ưu tiên chính xác $>$ chế tạo $>$ thẩm mỹ; đây
là lựa chọn thiết kế, không phải đại lượng học được. Để kiểm tra phép vô
hướng hóa này có tôn trọng cấu trúc đa mục tiêu hay không, chúng tôi tính
tương quan Spearman giữa điểm tổng hợp và tầng sắp xếp không trội
(non-dominated sorting) trên ba mục tiêu gốc. Phép tính được làm cho 24 mục
tiêu, mỗi mục tiêu 30 ứng viên (Hình~\ref{fig:spearman}). Cả 24 hệ số tương
quan đều dương (khoảng 0,43--0,89). Trung vị là 0,714, nhưng trung bình
0,693 thấp hơn một chút so với ngưỡng 0,7 đặt trước. Như vậy điểm tổng hợp
phù hợp ở mức vừa phải, chưa phải mạnh, với cấu trúc Pareto. Việc ứng viên
được chọn luôn nằm trên mặt Pareto thứ nhất là hệ quả toán học tất yếu với
trọng số dương, nên không được tính là bằng chứng.

\begin{table}[t]
\centering
\caption{Các chỉ báo đa mục tiêu. Khối trên: quy trình sản xuất có ép tuần
hoàn và chọn tốt nhất trong 30 mẫu, trên 300 mục tiêu kiểm tra. Tỉ lệ chế
tạo được: tỉ lệ mẫu sinh ra qua cả ba kiểm tra liên thông, kích thước tối
thiểu và tuần hoàn. Tỉ lệ trúng: tỉ lệ mục tiêu có mẫu đầu tiên là auxetic
($\nuot<0$). Khối dưới: mức phù hợp giữa điểm tổng hợp và tầng Pareto (24 mục
tiêu $\times$ 30 ứng viên).}
\label{tab:multi}
\small
\begin{tabular}{lccc}
\toprule
Mô hình & $\RR(\nuot)$, tốt nhất/30 & Tỉ lệ trúng mẫu đơn & Tỉ lệ chế tạo được\\
\midrule
cVAE sản xuất (bộ dữ liệu v2) & 0,995 & 0,997 & 0,247\\
cVAE trước đó (bộ dữ liệu v1) & 0,998 & 0,983 & 0,280\\
\midrule
\multicolumn{4}{p{0.85\linewidth}}{Điểm tổng hợp so với tầng Pareto: Spearman trung bình 0,693, trung vị 0,714, khoảng 0,43--0,89; 24/24 mục tiêu dương.}\\
\bottomrule
\end{tabular}
\end{table}

\begin{figure}[t]
\centering
\includegraphics[width=0.5\linewidth]{figures/fig_spearman_vi.pdf}
\caption{Phân bố trên 24 mục tiêu của tương quan Spearman giữa điểm tổng hợp
và tầng Pareto của 30 ứng viên (các mục tiêu: $|\Delta\nuot|$ kiểm chứng,
điểm chế tạo, điểm thẩm mỹ). Đường chấm: ngưỡng đặt trước cho giá trị trung
bình.}
\label{fig:spearman}
\end{figure}

\subsection{Tổng quát hóa ngoài phân phối}

\paragraph{Mục tiêu chấp nhận được.}
Trước khi thử ngoại suy, chúng tôi xác định mục tiêu nào thực sự tồn tại.
Với vật liệu trực hướng hai chiều, điều kiện xác định dương của ma trận độ
mềm đòi hỏi $\nuot\nuto<1$ \citep{lempriere1968poisson}. Do đó mục tiêu đối
xứng $\nuot^{*}=\nuto^{*}\le-1$ là không khả thi. Trong dữ liệu huấn luyện,
giá trị âm nhất $\nuot=-1{,}95$ đi kèm $\nuto=-0{,}04$, và tích $\nuot\nuto$
lớn nhất là 0,43. Điều này quan trọng cho việc đánh giá. Một thí nghiệm
``auxetic cực trị'' trước đây trong dự án đã dùng các mục tiêu như
$(-2{,}1;\,-2{,}1)$ và $(-2{,}2;\,-1{,}6)$, và cả sáu mục tiêu đều vi phạm
giới hạn này. Thất bại ở đó vì vậy không nói lên điều gì về khả năng ngoại
suy, và chúng tôi loại nó khỏi bài này. Các mục tiêu biên độ ngoài dải dưới
đây là dị hướng, với $\nuto^{*}=-0{,}05$.

\paragraph{Theo dấu: cVAE tổng quát hóa được, truy xuất thì không.}
Mọi mẫu huấn luyện đều auxetic ($-1{,}95\le\nuot\le-0{,}002$). Với sáu mục
tiêu không auxetic $0{,}1\le\nuot^{*}\le0{,}8$, truy xuất láng giềng gần nhất
từ thư viện 57\,216 mẫu tất yếu trả về mẫu gần 0 nhất ($\nuot=-0{,}002$), sai
dấu ở mọi mục tiêu. cVAE (tốt nhất trong 10 mẫu) cho $\nuot$ dương ở cả sáu
mục tiêu (Bảng~\ref{tab:signflip}, Hình~\ref{fig:signflip}), với MAE 0,19 so
với 0,40 của truy xuất. Tuy nhiên, biên độ đạt được bão hòa quanh 0,3: với
$\nuot^{*}=0{,}8$, cVAE chỉ đạt 0,30. Như vậy cVAE ngoại suy được
\emph{hướng} của tính chất nhưng không đạt đủ biên độ. Thí nghiệm này dùng
một checkpoint trước đó (bộ dữ liệu v1) và sáu mục tiêu, nên chỉ nên xem là
bằng chứng định tính.

\begin{table}[t]
\centering
\caption{Mục tiêu OOD đổi dấu (dữ liệu huấn luyện chỉ có $\nuot<0$). $\nuot$
kiểm chứng FE của truy xuất láng giềng gần nhất và của cVAE (mẫu đầu tiên và
tốt nhất trong 10). Cột cuối: tỉ lệ chế tạo được trong 10 mẫu cVAE.}
\label{tab:signflip}
\small
\begin{tabular}{cccccc}
\toprule
$(\nuot^{*};\nuto^{*})$ & Truy xuất & cVAE, mẫu đầu & cVAE, tốt nhất/10 & Chế tạo\\
\midrule
(0,1; 0,1) & $-0{,}002$ & 0,082 & 0,104 & 0,6\\
(0,3; 0,3) & $-0{,}002$ & 0,183 & 0,204 & 0,6\\
(0,5; 0,5) & $-0{,}002$ & 0,191 & 0,280 & 0,8\\
(0,8; 0,8) & $-0{,}002$ & 0,241 & 0,301 & 0,6\\
(0,2; 0,5) & $-0{,}002$ & 0,225 & 0,194 & 0,9\\
(0,5; 0,2) & $-0{,}002$ & 0,109 & 0,177 & 0,6\\
\midrule
MAE($\nuot$) & 0,402 & 0,237 & \textbf{0,191} & \\
\bottomrule
\end{tabular}
\end{table}

\begin{figure}[t]
\centering
\includegraphics[width=0.5\linewidth]{figures/fig_signflip_vi.pdf}
\caption{Tổng quát hóa theo dấu với các mục tiêu đối xứng không auxetic.
Truy xuất bị giới hạn bởi thư viện huấn luyện ($\nuot\le-0{,}002$). cVAE đạt
đúng dấu nhưng bão hòa về biên độ. Nét đứt: đường đồng nhất.}
\label{fig:signflip}
\end{figure}

\paragraph{Theo biên độ: mọi phương pháp bão hòa ở biên huấn luyện.}
Chúng tôi quét $\nuot^{*}\in\{-1{,}0;$ $-1{,}5;$ $-1{,}75;$ $-2{,}0;$
$-2{,}25;$ $-2{,}5\}$ ($\nuto^{*}=-0{,}05$), với 10 lần rút vector ẩn độc lập
mỗi mục tiêu, chọn tốt nhất trong 30 mẫu và ép tuần hoàn (Bảng~\ref{tab:ood},
Hình~\ref{fig:ood}). Trong dải huấn luyện, tinh chỉnh nhận biết nhị phân hóa
đưa trung vị sai số tương đối của $\nuot$ về 0,2--0,4\%. Vượt quá biên dưới
$-1{,}95$ của dữ liệu huấn luyện, mọi biến thể đều bão hòa. $\nuot$ kiểm
chứng âm nhất trong toàn bộ 30 thiết kế tại $\nuot^{*}=-2{,}5$ là $-1{,}99$.
Khi đó tinh chỉnh không kiểm soát còn kém hơn không tinh chỉnh (32,4\% so với
22,5\% tại $-2{,}5$). Biến thể có kiểm soát đạt ngưỡng 8\% đặt trước tại
$-2{,}0$ (4,9\%), nhưng chỉ vì mục tiêu này nằm cách biên 0,05. Từ $-2{,}25$
trở đi, nó chỉ ngang bằng bước chọn tốt nhất trong 30 mẫu. Tinh chỉnh là tìm
kiếm cục bộ trong không gian ẩn của bộ giải mã, nên không thể tạo ra biên độ
mà bộ giải mã chưa từng học để biểu diễn.

\begin{table}[t]
\centering
\caption{Quét biên độ dị hướng ($\nuto^{*}=-0{,}05$; tốt nhất trong 30 mẫu có
ép tuần hoàn; 10 lần lặp mỗi mục tiêu). Giá trị: trung vị sai số tương đối
của $\nuot$ kiểm chứng FE. Dải huấn luyện của $\nuot$ kết thúc tại $-1{,}95$.}
\label{tab:ood}
\small
\begin{tabular}{lcccccc}
\toprule
$\nuot^{*}$ & $-1{,}0$ & $-1{,}5$ & $-1{,}75$ & $-2{,}0$ & $-2{,}25$ & $-2{,}5$\\
\midrule
Tốt nhất/30, không tinh chỉnh         & 0,5\% & 0,9\% & 4,1\% & 11,0\% & 16,1\% & 22,5\%\\
+ tinh chỉnh liên tục                 & 1,4\% & 1,3\% & 3,2\% & 36,5\% & 38,9\% & 43,9\%\\
+ tinh chỉnh nhận biết nhị phân hóa   & 0,4\% & 0,2\% & 0,3\% & 12,5\% & 22,3\% & 32,4\%\\
+ nhận biết nhị phân hóa, có kiểm soát & 0,2\% & 0,2\% & 0,3\% & \textbf{4,9\%} & 14,8\% & 21,9\%\\
\bottomrule
\end{tabular}
\end{table}

\begin{figure}[t]
\centering
\includegraphics[width=0.62\linewidth]{figures/fig_ood_vi.pdf}
\caption{Trung vị sai số tương đối của $\nuot$ kiểm chứng FE theo mục tiêu
($\nuto^{*}=-0{,}05$). Thanh sai số: khoảng tứ phân vị trên 10 lần rút vector
ẩn mỗi mục tiêu; các điểm được dịch ngang để dễ đọc. Vùng tô: ngoài dải huấn
luyện. Đường chấm: ngưỡng 8\% đặt trước tại $\nuot^{*}=-2{,}0$.}
\label{fig:ood}
\end{figure}

\subsection{Thí nghiệm loại bỏ về kiến trúc}

\paragraph{Đầu KAN so với đầu tuyến tính.}
Chúng tôi thay ba lớp tuyến tính của cVAE (trung bình ẩn, log-phương sai và
nút cổ chai của bộ giải mã) bằng các lớp Kolmogorov--Arnold
\citep{liu2024kan}. Cả hai loại đầu được huấn luyện với cùng công thức, siêu
tham số và seed (123 và 7). KAN chỉ được coi là thắng nếu KTC ghép cặp không
chứa 0 ở cả hai seed. Kết quả là KAN không thắng ở chế độ nào
(Bảng~\ref{tab:kan}). Ở chế độ sản xuất (tốt nhất trong 30 mẫu, không tinh
chỉnh), đầu tuyến tính tốt hơn có ý nghĩa ở cả hai seed. Khi có tinh chỉnh,
cả hai đều đạt $\RR(\nuot)=0{,}997$--$0{,}999$. Đầu KAN có số tham số gấp 3,2
lần (4,79\,triệu so với 1,49\,triệu) và thời gian mỗi epoch gấp khoảng 1,6
lần, nên chúng tôi giữ đầu tuyến tính.

\begin{table}[t]
\centering
\caption{Thí nghiệm loại bỏ KAN so với tuyến tính trên IN100 (hai seed).
$\Delta$MAE: mức giảm tương đối MAE($\nuot$) của KAN so với tuyến tính (dương
là KAN tốt hơn), KTC 95\% bootstrap ghép cặp. Chữ đậm: KTC không chứa 0.}
\label{tab:kan}
\small
\setlength{\tabcolsep}{4pt}
\begin{tabular}{lcccc}
\toprule
 & \multicolumn{2}{c}{$\RR(\nuot)$ Tuyến tính / KAN} & \multicolumn{2}{c}{$\Delta$MAE KAN so với tuyến tính}\\
\cmidrule(lr){2-3}\cmidrule(lr){4-5}
Chế độ & seed 123 & seed 7 & seed 123 & seed 7\\
\midrule
Mẫu đơn lẻ                 & 0,852 / 0,904 & 0,815 / 0,911 & $+13{,}3\%$ [$-0{,}2$; 24,4] & $\mathbf{+18{,}2\%}$ [3,8; 30,4]\\
Mẫu đơn lẻ + tinh chỉnh    & 0,994 / 0,997 & 0,998 / 0,999 & $+33{,}2\%$ [$-15{,}1$; 60,5] & $\mathbf{+33{,}9\%}$ [11,7; 52,3]\\
Tốt nhất/30                & 0,977 / 0,974 & 0,982 / 0,980 & $\mathbf{-25{,}2\%}$ [$-56{,}2$; $-1{,}0$] & $\mathbf{-19{,}9\%}$ [$-42{,}1$; $-1{,}7$]\\
Tốt nhất/30 + tinh chỉnh   & 0,998 / 0,997 & 0,998 / 0,999 & $+11{,}3\%$ [$-38{,}7$; 43,6] & $\mathbf{+44{,}1\%}$ [24,9; 59,7]\\
\bottomrule
\end{tabular}
\end{table}

\paragraph{Bộ giải mã biểu diễn ẩn wavelet.}
Một bộ giải mã WIRE không phụ thuộc độ phân giải \citep{saragadam2023wire}
được thử lần cuối. Lần thử này kết hợp phép chiếu Heaviside trong loss vật
lý, không ép đối xứng và công thức hai giai đoạn. Tiêu chí đặt trước là tỉ
lệ trúng mẫu đơn tối thiểu 30\% và $\RR$ FE dương trên 24 mục tiêu kiểm tra.
Kết quả đạt tỉ lệ trúng 12,5\% và $\RR$ FE bằng $-7{,}19$ sau bước chọn tốt
nhất trong 30 mẫu, trong khi bộ giải mã tích chập đạt tỉ lệ trúng khoảng
99,7\%. Chỉ 4,6\% mẫu của nó chế tạo được. Vì vậy hướng này được dừng lại.

% ---------------------------------------------------------------------------
\section{Thảo luận}

\paragraph{Độ chính xác.}
Phát hiện phương pháp luận trung tâm là: khi tinh chỉnh một thiết kế được
sinh ra dựa trên vật lý, hàm mục tiêu phải được tính trên đúng đối tượng sẽ
được kiểm chứng. Lấy đạo hàm qua bộ giải chính xác loại bỏ được sai số
surrogate, nhưng không loại bỏ được sự lệch giữa đầu ra liên tục của bộ giải
mã và ô cơ sở đã nhị phân hóa. Hình~\ref{fig:gap} cho thấy sự lệch này cho
phép bộ tối ưu đạt hàm mục tiêu $10^{-10}$ trên một trường mà bộ kiểm chứng
nhìn thấy rất khác. Vấn đề không lộ ra trên đường cong loss, và chỉ bộc lộ
khi tinh chỉnh một baseline mạnh. Các thành phần của cách khắc phục là tiêu
chuẩn trong tối ưu hóa tô-pô dựa trên mật độ
\citep{guest2004achieving,wang2011projection}. Điểm mới là áp dụng chúng lên
đầu ra của một mô hình sinh cố định trong quá trình tìm kiếm không gian ẩn,
cùng với đúng phép lấy mẫu lại và hậu xử lý tuần hoàn dùng khi kiểm chứng.
Nguyên tắc này hẳn cũng mở rộng được cho bất kỳ bước hậu xử lý nào nằm trong
khâu kiểm chứng.

\paragraph{Đa mục tiêu.}
Quy trình xử lý ba mục tiêu ở các giai đoạn khác nhau. Tính tuần hoàn được
đảm bảo bằng cấu trúc, tính liên thông và kích thước chi tiết đi vào qua
điểm tổng hợp, còn độ chính xác được cải thiện sau cùng bằng tinh chỉnh vẫn
giữ biên tuần hoàn. Cách sắp xếp này mang tính thực dụng, và bằng chứng của
chúng tôi cho nó mới là một phần. Điểm tổng hợp chỉ phù hợp vừa phải với xếp
hạng Pareto. Một nguyên nhân có thể là việc chuẩn hóa theo thứ hạng của số
hạng độ chính xác trong từng tập ứng viên, nhưng chúng tôi chưa kiểm chứng
điều này. Chúng tôi cũng chưa đo tinh chỉnh có làm thay đổi tính liên thông
hay kích thước chi tiết hay không. Người thiết kế cần đảm bảo chắc chắn khả
năng chế tạo vẫn có thể giới hạn bước chọn vào các ứng viên qua mọi kiểm
tra. Trong các thí nghiệm trước, cách này cần tập ứng viên lớn
($N\approx1500$) và làm giảm độ chính xác.

\paragraph{Tổng quát hóa ngoài phân phối.}
Kết quả của chúng tôi gợi ý một phát biểu chính xác hơn cho nhận định phổ
biến rằng mô hình sinh ``tổng quát hóa được''. Theo dấu của hệ số Poisson,
cVAE tạo ra thiết kế trong vùng có mật độ dữ liệu huấn luyện bằng 0, điều mà
phương pháp tra cứu không thể làm được về mặt cấu trúc. Theo biên độ, mô
hình sinh, bước tinh chỉnh và bước chọn tốt nhất trong $N$ mẫu đều dừng ở
biên dữ liệu huấn luyện. Vì vậy các tuyên bố về OOD nên nêu rõ hướng trong
không gian tính chất, và nên kiểm tra trước rằng mục tiêu là chấp nhận được
về vật lý. Ngoại suy biên độ nhiều khả năng cần dữ liệu huấn luyện mới gần
biên, ví dụ từ tối ưu hóa tô-pô khởi tạo bằng các thiết kế được sinh ra,
hơn là lấy mẫu tốt hơn từ mô hình hiện tại.

\paragraph{Giới hạn.}
(i) Mọi phân tích dùng đồng nhất hóa đàn hồi tuyến tính, biến dạng nhỏ, cho
ô cơ sở hai chiều trên lưới $50\times50$ với hệ số Poisson vật liệu nền 0,3.
Biến dạng lớn và kiến trúc ba chiều chưa được xét.
(ii) Thiết kế được kiểm chứng bằng số, kể cả với một bản cài đặt FE độc lập,
nhưng chưa được kiểm chứng thực nghiệm.
(iii) Thí nghiệm đổi dấu dùng sáu mục tiêu và một checkpoint trước đó; phép
quét biên độ dùng 10 lần lặp mỗi mục tiêu.
(iv) Điểm thẩm mỹ là một heuristic đơn giản (đối xứng gương và độ trơn biên),
chưa được đối chiếu với đánh giá của con người.
(v) Trong phân phối, truy xuất láng giềng gần nhất rất cạnh tranh về độ
chính xác, nên lợi thế của mô hình sinh thể hiện ở đây chỉ giới hạn ở việc
đổi dấu ngoài phân phối.
(vi) Tinh chỉnh tốn thêm khoảng 77 lần giải FE mỗi mục tiêu. Chi phí này rẻ
ở 2D nhưng sẽ đòi hỏi bộ giải adjoint hiệu quả ở 3D.

% ---------------------------------------------------------------------------
\section{Phương pháp}

\subsection{Sinh dữ liệu}
Ô cơ sở được sinh bằng tối ưu hóa tô-pô SIMP
\citep{bendsoe1989optimal,sigmund2001code,andreassen2011efficient} với bộ lọc
mật độ \citep{bourdin2001filters}, điều kiện biên tuần hoàn và đồng nhất hóa
dựa trên năng lượng
\citep{guedes1990preprocessing,andreassen2014homogenization,xia2015design}.
Các hàm mục tiêu auxetic được cực tiểu hóa từ bốn tô-pô khởi tạo, dùng tiêu
chuẩn tối ưu (OC) hoặc phương pháp tiệm cận di động \citep{svanberg1987mma}
tùy tô-pô. Các tham số quy trình (tỉ lệ thể tích, hệ số phạt, bán kính lọc,
giới hạn bước, kích thước lỗ rỗng, góc xoay mạng) được sàng lọc bằng lấy mẫu
siêu lập phương Latin \citep{mckay1979lhs}. Sau đó chúng được tinh chỉnh
trong một quy trình thiết kế thí nghiệm thích nghi nhiều lô, dùng dãy Sobol'
\citep{sobol1967distribution} và siêu lập phương Latin tối ưu. Các mẫu hội
tụ dao động, tô-pô suy biến hoặc nhãn ngoài dải vật lý bị loại bỏ. Trường
mật độ $50\times50$ được lọc hộp về ảnh $64\times64$ và gán nhãn
$(\nuot,\nuto)$ đồng nhất hóa. Bộ dữ liệu sản xuất gồm 13\,624 ô cơ sở phân
biệt, chia 70/15/15 theo họ tô-pô khởi tạo (2\,044 mẫu kiểm tra). Ảnh huấn
luyện được tăng cường sáu lần bằng phép xoay và đối xứng; phép xoay
$90^{\circ}$ hoán đổi $\nuot\leftrightarrow\nuto$. Kết quả là 57\,216 mẫu
huấn luyện.

\subsection{Đồng nhất hóa và độ nhạy chính xác}
Trường mật độ $\rho\in[0,1]^{50\times50}$ được rời rạc hóa bằng phần tử tứ
giác song tuyến tính với nội suy SIMP cải tiến
$E(\rho)=E_{\min}+\rho^{p}(E_0-E_{\min})$, trong đó $E_{\min}=10^{-9}E_0$,
$p=3$ và $\nu_0=0{,}3$. Tensor độ cứng đồng nhất hóa $\mathbf{Q}$ nhận được
từ ba trường hợp tải tuần hoàn thông qua năng lượng tương hỗ phần tử
\citep{xia2015design}. Hệ số Poisson lấy từ $\mathbf{S}=\mathbf{Q}^{-1}$:
$\nuot=-S_{12}/S_{11}$ và $\nuto=-S_{12}/S_{22}$. Công thức này vẫn đúng với ô
cơ sở bị xoay, dị hướng hoàn toàn. Nhờ công thức năng lượng tự liên hợp, độ
nhạy $\partial\mathbf{Q}/\partial\rho_e$ có sẵn mà không cần giải thêm. Độ
nhạy được lan truyền giải tích qua
$\partial\mathbf{S}=-\mathbf{S}\,(\partial\mathbf{Q})\,\mathbf{S}$ và quy tắc
đạo hàm thương, trong một hàm autograd tùy biến của PyTorch
\citep{paszke2019pytorch}. Gradient của hàm này khớp sai phân hữu hạn với sai
số tương đối dưới $10^{-4}$. Bộ giải khớp với một bản cài đặt độc lập bằng
scikit-fem \citep{gustafsson2020skfem} tới $10^{-9}$ trên 24 thiết kế.

\subsection{cVAE}
cVAE \citep{kingma2014vae,sohn2015cvae} được điều kiện hóa theo
$c=(\nuot,\nuto)$ và có không gian ẩn 32 chiều. Bộ mã hóa gồm bốn khối tích
chập--chuẩn hóa batch--ReLU--pooling (32 đến 256 kênh), giữ lại bản đồ không
gian $4\times4$, tiếp theo là các lớp tuyến tính cho trung bình và
log-phương sai. Bộ giải mã đối xứng với bộ mã hóa, dùng tích chập chuyển vị
và đầu ra sigmoid. Giai đoạn 1 huấn luyện với loss tái tạo entropy chéo nhị
phân, số hạng KL có khởi động tuyến tính \citep{bowman2016generating}, và số
hạng nhất quán tính chất từ một surrogate tích chập cố định (trọng số
$\gamma=20$; surrogate đạt $\RR$ kiểm tra 0,974 cho $\nuot$ và 0,964 cho
$\nuto$). Giai đoạn 2 tiếp tục từ giai đoạn 1 và thêm bình phương sai số
giữa $c$ và hệ số Poisson đồng nhất hóa FE của các mẫu \emph{tiên nghiệm}
$z\sim\mathcal N(0,I)$ đã giải mã (trọng số 20; 8 mẫu mỗi hai epoch). Đây
chính là chế độ dùng khi suy luận, và huấn luyện loss này từ đầu thì thất
bại. Tối ưu dùng Adam \citep{kingma2015adam}. Checkpoint được chọn theo $\RR$
kiểm chứng FE trên các điều kiện kiểm định chứ không theo loss kiểm định, vì
khởi động KL làm loss kiểm định không đáng tin.

\subsection{Kiểm chứng, khả năng chế tạo và điểm tổng hợp}
Một ảnh giải mã được kiểm chứng qua bốn bước: (1) ép biên tuần hoàn, thay các
cột rồi các hàng biên đối diện bằng trung bình của chúng; (2) lấy ngưỡng
0,5; (3) lấy mẫu lại láng giềng gần nhất về $50\times50$ bằng bảng chỉ số của
thư viện ảnh dùng khi xây dựng bộ dữ liệu; và (4) đồng nhất hóa FE. Khả năng
chế tạo kết hợp ba kiểm tra trên ảnh nhị phân: pha rắn liên thông duy nhất,
kích thước chi tiết tối thiểu hai pixel (phép co hình thái), và tính tuần
hoàn (tối đa 10\% pixel lệch giữa các cạnh đối diện). Điểm có bậc
$s_{\mathrm{man}}$ là trung bình của ba chỉ báo. Điểm thẩm mỹ
$s_{\mathrm{aes}}$ lấy trung bình giữa độ đối xứng gương theo hai trục và
nghịch đảo tỉ lệ chu vi/diện tích đã chuẩn hóa, tính với biên tuần hoàn. Điểm
chính xác $s_{\mathrm{acc}}$ là $|\Delta\nuot|$ chuẩn hóa min--max trong tập
ứng viên, đổi dấu để giá trị cao hơn là tốt hơn. Phép kiểm tra Pareto áp
dụng sắp xếp không trội lặp trên ba mục tiêu gốc, rồi tính tương quan giữa
tầng thu được và điểm tổng hợp cho từng mục tiêu.

\subsection{Tinh chỉnh không gian ẩn}
Cho $\boldsymbol\nu^{*}$ và vector ẩn xuất phát $z_0$ (mẫu đơn lẻ hoặc ứng
viên tốt nhất trong $N$ mẫu), chúng tôi cực tiểu hóa
\begin{equation}
\mathcal L_\beta(z)=\tfrac12\bigl\lVert
\boldsymbol\nu\bigl(\mathcal R\circ H_\beta\circ\mathcal P\circ D(z,c)\bigr)
-\boldsymbol\nu^{*}\bigr\rVert_2^2 ,
\end{equation}
trong đó $D$ là bộ giải mã cố định, $\mathcal P$ là phép ép biên tuần hoàn
(tuyến tính), và $\mathcal R$ là phép lấy mẫu lại láng giềng gần nhất được
cài đặt như một phép gather khả vi với cùng bảng chỉ số. Cuối cùng,
\begin{equation}
H_\beta(\rho)=\frac{\tanh(\beta\eta)+\tanh\bigl(\beta(\rho-\eta)\bigr)}
{\tanh(\beta\eta)+\tanh\bigl(\beta(1-\eta)\bigr)},\qquad \eta=0{,}5,
\end{equation}
là phép chiếu Heaviside làm trơn \citep{wang2011projection}. 32 bước L-BFGS
(bước 0,5, tìm kiếm đường strong-Wolfe) được chia đều cho
$\beta\in\{1;4;16;64\}$. L-BFGS được khởi động lại ở mỗi mức, vì thông tin độ
cong từ phép chiếu mềm hơn không còn đúng cho mức mới. Baseline liên tục áp
loss lên $D(z,c)$ với nội suy song tuyến tính trong 30 bước. Tinh chỉnh có
kiểm soát chỉ giữ thiết kế đã tinh chỉnh nếu $|\Delta\nuot|+|\Delta\nuto|$
kiểm chứng của nó nhỏ hơn thiết kế xuất phát.

\subsection{Giao thức ngoài phân phối và baseline truy xuất}
Phương pháp truy xuất trả về mẫu huấn luyện có khoảng cách Euclid nhỏ nhất
tới mục tiêu trong không gian $(\nuot,\nuto)$. Thí nghiệm đổi dấu dùng sáu
mục tiêu không auxetic, 10 mẫu cVAE mỗi mục tiêu và chọn bằng FE. Phép quét
biên độ dùng sáu mục tiêu dị hướng với $\nuto^{*}=-0{,}05$ và 10 lần lặp độc
lập của bước chọn tốt nhất trong 30 mẫu cho mỗi mục tiêu. Tính chấp nhận
được của mục tiêu được kiểm tra theo điều kiện $\nuot\nuto<1$
\citep{lempriere1968poisson}.

\subsection{Thống kê và đăng ký trước tiêu chí}
IN100 gồm 100 điều kiện kiểm tra rút với seed 123. Mức giảm MAE tương đối và
KTC 95\% được tính bằng bootstrap phân vị ghép cặp theo điều kiện, với
10\,000 lần lấy mẫu lại \citep{efron1993bootstrap}. Một cải thiện chỉ được
khẳng định khi KTC không chứa 0. Các tiêu chí được ghi lại trước mỗi thí
nghiệm, gồm: giảm MAE ít nhất 20\% với mẫu đơn lẻ; MAE không tăng có ý nghĩa
sau bước chọn tốt nhất trong 30 mẫu; trung vị sai số tương đối tối đa 8\%
tại $\nuot^{*}=-2{,}0$; tương quan Spearman trung bình tối thiểu 0,7 cho điểm
tổng hợp; với KAN, KTC không chứa 0 ở cả hai seed; với WIRE, tỉ lệ trúng tối
thiểu 30\% và $\RR$ FE $>0$. L-BFGS trên GPU không tất định từng bit, và các
lần chạy lặp lệch nhau khoảng 0,006 về $\RR$.

\section*{Dữ liệu và mã nguồn}
Mã nguồn, checkpoint đã huấn luyện và toàn bộ kết quả theo từng điều kiện
dùng cho các bảng và hình sẽ được công bố tại [đường dẫn kho mã - sẽ bổ
sung].

\section*{Lời cảm ơn và xung đột lợi ích}
[Sẽ bổ sung.]

\bibliographystyle{unsrtnat}
\bibliography{refs}

\end{document}
